% titel: Blatt 9, A5: Kc, Kp, kleinstes Zwang, dynamisches Gleichgewicht, Q
% skript: Mortimer Kap. 14 (chemisches Gleichgewicht)

% Teil I: Werkzeuge, kein Lösungsweg.

Die Grössen rund um das chemische Gleichgewicht präzise abzugrenzen.

\begin{werkzeugkasten}[Begriffe des chemischen Gleichgewichts]
Für eine allgemeine Reaktion $aA + bB \rightleftharpoons cC + dD$:
\begin{itemize}
  \item \textbf{K$_c$:} Gleichgewichtskonstante bezogen auf
        Konzentrationen
        \[
          K_c = \frac{[C]^c[D]^d}{[A]^a[B]^b}
        \]
        (Gleichgewichtskonzentrationen).
  \item \textbf{K$_p$:} Gleichgewichtskonstante bezogen auf Partialdrücke
        ($K_p = \frac{p_C^c p_D^d}{p_A^a p_B^b}$).
  \item \textbf{Prinzip des kleinsten Zwanges (Le Chatelier):} ein von aussen
        aufgezwungener Zwang (Änderung von $p$, $T$, $c$, ...) verschiebt
        das Gleichgewicht so, dass der Zwang abgeschwächt wird.
  \item \textbf{Dynamisches Gleichgewicht:} Hin- und Rückreaktion laufen
        mit gleicher Geschwindigkeit; die Nettoänderung ist null, aber es
        reagiert ständig weiter.
  \item \textbf{Reaktionsquotient Q:} wie $K_c$, aber mit den
        \emph{aktuellen} (nicht Gleichgewichts-)Konzentrationen; er zeigt
        die Richtung an ($Q<K \Rightarrow$ Reaktion in Richtung Produkte).
\end{itemize}
\textit{Skript:} Mortimer Kap. 14.
\end{werkzeugkasten}

\begin{achtung}[Darauf achten]
\begin{itemize}
  \item Nur $T$-Änderungen verändern $K_c$/$K_p$ selbst; $c$- oder
        $p$-Änderungen verschieben nur das Gleichgewicht, $K$ bleibt.
  \item Q vergleicht mit $K$: $Q<K$ vorwärts, $Q>K$ rückwärts, $Q=K$
        Gleichgewicht.
  \item Katalysator beschleunigt Hin- und Rückreaktion gleich --- kein
        Einfluss auf die Lage.
\end{itemize}
\end{achtung}
